\scanpage{58}
% @01–02, category
Soit $(SS)$ la catégorie des ensembles semi-simpliciaux
% @03, presheaves
\[
 (SS)=\widehat{\mathbb{\Delta}}=\underline{\operatorname{Hom}}(\mathbb{\Delta}^{\circ},(\mathrm{Ens}))
\]
% @04, topos
C'est un topos, et si
% @05, object
\[
 \unsure{\mathscr{X}|\mathcal{X}|X}\in\operatorname{Ob}(SS)
\]
% @06–07, slice
est un ens $1/2$ simplicial, on regarde le topos induit
% @07, slice identification
\[
 (SS)/\mathscr{X}\simeq\widehat{\mathbb{\Delta}/\mathscr{X}}\simeq
 \underline{\operatorname{Hom}}(\mathbb{\Delta}^\circ,\mathrm{Ens})
\]
% @08–17, sheaves on the object
Pour nous, l'étude cohomologique de $\mathscr{X}$ est celle de ce topos induit. Les faisceaux \inserted{(d'ens)} ``sur $\mathscr{X}$'' sont les \deleted{\unsure{gr}} faisceaux semi-simpliciaux sur $\mathscr{F}\to\mathscr{X}$ sur $\mathscr{X}$, ou encore les foncteurs $(\mathbb{\Delta}/\mathscr{X})^\circ\to(\mathrm{Ens})$ ; les faisceaux abéliens sur $\mathscr{X}$ sont les objets groupes abéliens $F\to\mathscr{X}$ sur $\mathscr{X}$, ou encore les foncteurs $(\mathbb{\Delta}/\mathscr{X})^\circ\to(\mathrm{Ab})$ etc. Idem pour faisceaux abéliens gradués différentiels etc. Si $A_{\deleted{\unsure{*|\bullet}}}$ est un ens. ssl. on peut regarder
% @18, product
\[
 A\times\mathscr{X}\longrightarrow\mathscr{X}\quad\text{comme faisceau sur }\mathscr{X}
\]
% @19–21, induced sheaf
(c'est aussi le faisceau induit par $A_{\unsure{\alpha|*|\bullet}}$ sur le topos induit $(SS/\mathscr{X})$), c'est aussi le foncteur composé
% @22, composite
\[
 (\mathbb{\Delta}^{\circ}/\mathscr{X})^\circ\longrightarrow
 (\mathbb{\Delta})^\circ\xrightarrow{A}(\mathrm{Ens}).
\]
% archive sheet number: 57

